% toolwall-arxiv.tex
% Toolwall: A Defense-in-Depth Library for Indirect Prompt Injection
% at the Tool-Response Boundary
%
% arXiv: cs.CR (primary) + cs.AI (cross-list)
% Target length: ~8 pages including references (NeurIPS / ACL workshop format)
%
% Build:
%   cd paper/
%   pdflatex toolwall-arxiv.tex && bibtex toolwall-arxiv && \
%   pdflatex toolwall-arxiv.tex && pdflatex toolwall-arxiv.tex
%
% The .bib is inline below \begin{thebibliography} to keep the
% preprint self-contained for first arXiv submission. A CITATION.cff
% at repo root mirrors the citation block.

\documentclass[11pt,letterpaper]{article}

\usepackage[margin=1in]{geometry}
\usepackage[utf8]{inputenc}
\usepackage[T1]{fontenc}
\usepackage{microtype}
\usepackage{lmodern}
\usepackage{booktabs}
\usepackage{array}
\usepackage{multirow}
\usepackage{xcolor}
\usepackage{listings}
\usepackage[hidelinks]{hyperref}
\usepackage{url}
\usepackage{enumitem}
\usepackage{authblk}
\usepackage{abstract}
\usepackage{titlesec}

\titleformat{\section}{\normalfont\Large\bfseries}{\thesection}{0.6em}{}
\titleformat{\subsection}{\normalfont\large\bfseries}{\thesubsection}{0.6em}{}

\lstset{
  basicstyle=\ttfamily\footnotesize,
  breaklines=true,
  columns=fullflexible,
  keepspaces=true,
  showstringspaces=false,
  frame=single,
  framerule=0.3pt,
  rulecolor=\color{black!40},
  framesep=4pt,
  xleftmargin=4pt,
  xrightmargin=4pt,
  aboveskip=6pt,
  belowskip=6pt,
}

\title{\bfseries Toolwall: A Defense-in-Depth Library for Indirect Prompt
Injection at the Tool--Response Boundary}

\author[1,2]{Muhammad Hasnain}
\affil[1]{bryteai.studio}
\affil[2]{Bin Abdullah LLC, Brooklyn, NY \ \ \textemdash\ \ NVIDIA Inception member}
\affil[ ]{\texttt{hasnain@bryteai.studio}}

\date{Preprint --- \today}

\begin{document}

\maketitle

\begin{abstract}
We present \textsc{toolwall}, the first opensource (Apache-2.0)
indirect-prompt-injection defense library positioned at the
agent$\leftrightarrow$tool response boundary --- the surface that
EchoLeak (CVE-2025-32711, CVSS~9.3), MCP tool poisoning,
full-schema poisoning, rug-pull description mutation, and
markdown-image exfiltration all exploit, and that existing
opensource defenses (Llama-Guard-3, Prompt-Guard-86M,
Spotlighting~\cite{hines2024spotlight}, NeMo
Guardrails~\cite{nemoguard2024}) either do not cover or cover
only at the input/output side of the model. A single decorator
(\verb|@toolwall(policy=...)|) yields three layers --- policy,
IPI detector, enforcement --- with a four-line integration
footprint, framework-agnostic adapters for the Claude Agent SDK,
CrewAI, LangGraph, OpenAI Agents SDK, and the Model Context
Protocol (MCP), and zero modifications to the host agent.

On the 200-case in-house split of our frozen
\textsc{toolwall-bench-v0.1} (185~attacks + 15~benigns, seven
attack families), \textsc{toolwall} (rules + LoRA judge) achieves
attack success rate ASR $= 0.097$ with Wilson $95\%$ CI
$[0.062, 0.149]$ at utility $0.933$ $[0.702, 0.988]$. The best
publicly available OSS baseline we could run on the same split,
\texttt{Niansuh/Prompt-Guard-86M}, reaches ASR $0.427$ $[0.358,
0.499]$ at utility $0.467$ $[0.248, 0.699]$: \textsc{toolwall}'s
upper CI ($0.149$) is strictly below PromptGuard's lower CI
($0.358$), so the gap is non-overlapping at $n=185$. The rule
layer alone is a measured floor of ASR $0.778$ $[0.713, 0.832]$,
showing that the judge layer carries the attack-class coverage
pattern matching cannot.

We additionally report a MinHash + 5-gram Jaccard audit of the
train$\leftrightarrow$test corpora (overall mean $0.006$, max
$0.447$, only $2/200$ cases above the leak threshold), and an
honest out-of-distribution probe of $10$ CrewAI-generated attacks
from $10$ novel attack families absent from training, where the
rule-only floor degrades to ASR $0.80$ $[0.490, 0.943]$. The
dev-split saturation of the judge LoRA ($1.0/1.0/1.0$) is
documented and treated as suspect; the $200$-case test split is
the only authoritative number. Code, weights, dataset,
leaderboard, and a deterministic Playwright recorder are public
so the community can attack and improve the layer.

\noindent\textbf{Code:} \url{https://github.com/bryteai/toolwall}
\quad
\textbf{Demo:} \url{https://fence.bryteai.studio}
\end{abstract}

\section{Introduction}
\label{sec:intro}

LLM agents that call external tools (HTTP fetchers, file readers,
MCP servers, database query tools) inherit the trust boundary of
\emph{whatever string the tool returns}. When that string contains
an attacker-authored instruction, the agent has no native mechanism
to refuse: alignment training conditions the model to follow plain
text, and the tool response \emph{is} plain text. This is indirect
prompt injection (IPI), and 2025--2026 has produced an unusually
concrete evidence pile that it is no longer hypothetical:

\begin{itemize}[leftmargin=*,nosep]
\item \textbf{EchoLeak} (CVE-2025-32711, CVSS~9.3,
patched June~2025): zero-click data exfiltration from M365 Copilot
by embedding an instruction in an untrusted email that caused
auto-fetched markdown images to leak calendar/email contents.
First named production IPI exploit.
\item \textbf{CVE-2026-30615}: MCP supply-chain RCE chain
affecting Cursor, VS~Code, Windsurf, Claude~Code, and Gemini-CLI;
OX~Security disclosed $\sim$200K vulnerable MCP server instances.
\item \textbf{MCPTox} \cite{mcptox2025}: 353 real MCP tools $\times$
1{,}312 attacks. \texttt{o1-mini} ASR $72.8\%$;
\texttt{Claude-3.7-Sonnet} refused only $2.7\%$ of attacks.
\item \textbf{OWASP ASI Top~10 (Q1~2026)} \cite{owaspasi2026}:
A new framework that elevates IPI, Tool Poisoning, and Memory
Tampering to positions A01--A03.
\item \textbf{Anthropic MCP STDIO RCE class} (April~2026): 10~CVEs
across LiteLLM, LangChain, LangFlow, Flowise, LettaAI, and
LangBot. Anthropic declined a protocol-level fix
(``expected behaviour''), confirming the gap is a library
opportunity, not a vendor one.
\end{itemize}

The defensive landscape today is bifurcated. \emph{Closed} platforms
(Lakera, Pillar, Invariant, Northflank) require integration with a
SaaS endpoint and do not publish their patterns or model weights.
\emph{Open} tooling is either upstream of the response
(MCP-scan-style metadata audit) or runs inside the model
(prompt-sandwiching, Spotlighting~\cite{hines2024spotlight}). No
Apache-2.0 library currently \emph{inspects the response body} at
the tool-call boundary before it re-enters the model's context.

\textsc{toolwall} is that primitive. Our contributions:
\begin{enumerate}[leftmargin=*,nosep]
\item A three-layer architecture (policy / detector / enforcement)
exposed via a single decorator with a four-line integration
footprint (\S\ref{sec:design}).
\item Reference implementations in Python and TypeScript with
adapters for the Claude Agent SDK, CrewAI, LangGraph, OpenAI~Agents
SDK, FastMCP (Python), \texttt{@modelcontextprotocol/sdk}
(TypeScript), and a Cloudflare Workers MCP-server scaffold
(\S\ref{sec:impl}).
\item \textsc{toolwall-bench-v0.1}: 1{,}054 IPI cases
($500$ InjecAgent~\cite{zhan2024injecagent} +
$354$ AgentDojo~\cite{debenedetti2024agentdojo} +
$200$ new MCP-specific cases), frozen 600/200/254 train/dev/test
splits with SHA256-pinned test (\S\ref{sec:eval}).
\item Transparent rule-only results, a deterministic Playwright
video recorder for reproducible demos, and an Apache-2.0 release
with no proprietary patterns (\S\ref{sec:eval}--\ref{sec:limits}).
\end{enumerate}

\section{Threat model}
\label{sec:threat}

The full threat model is shipped as \texttt{THREAT\_MODEL.md} in
the repository and was adversarially reviewed at v0.1.0 cut. The
operative shape is:

\paragraph{Attacker capabilities.}
The adversary controls the \emph{content} of any tool response
returned to the agent: web pages, RAG documents, email bodies,
API JSON, file bodies, MCP tool output. The adversary is assumed
to know the agent's system prompt and approximate tool schema
(worst case). The adversary \emph{cannot} modify model weights,
the agent's system prompt, or the agent's code.

\paragraph{Defender position.}
\textsc{toolwall} sits between every tool invocation and the model
context window. It sees raw tool output \emph{before} it is
serialized back to the LLM. It owns policy
(allowlists, scopes, budgets, redactions) and the IPI detector
that runs on each response.

\paragraph{In-scope attacks.}
\begin{enumerate}[leftmargin=*,nosep]
\item Indirect prompt injection in tool responses (page bodies,
file contents, JSON values).
\item Tool poisoning / rug-pull --- the server changes a tool
description post-approval~\cite{invariant2025toolpoisoning}.
\item Full-schema poisoning --- malicious schema fields embedded
in tool metadata~\cite{cyberark2025schema}.
\item Exfiltration via tool side effects (POST to attacker URL,
markdown-image GET, secret encoding into search queries).
\item Resource abuse --- runaway tool loops, unbounded
egress / token cost.
\item Capability escalation through tool-call shape violations.
\end{enumerate}

\paragraph{Explicit non-goals.}
We do \emph{not} defend against (i)~direct adversarial-suffix
jailbreaks of the underlying LLM (HarmBench /
JailbreakBench~\cite{harmbench2024,jbb2024} territory),
(ii)~host-OS or sandbox escapes (Codex CLI \texttt{notify}-config
class~\cite{cymulate2026codex}),
(iii)~authentication / IAM,
(iv)~network packet filtering. \textsc{toolwall} is an
\emph{application-layer} firewall on tool semantics; it must be
deployed with a sandbox, an authenticated tool layer, and a
reasonably aligned model. None of those alone is sufficient and
neither is \textsc{toolwall} alone.

\paragraph{Where \textsc{toolwall} fails open.}
We list, by name, the failure modes a Stanford-tier reviewer
should ask us about. (1)~Judge-model evasion: a sufficiently
small / paraphrased instruction may slip past both rules and
judge; the audit log will retain evidence but the response will
pass. (2)~Encrypted payloads: we redact base64 but not arbitrary
ciphertext. (3)~Multi-hop attacks where the malicious instruction
is split across two tool responses are detected only if either
fragment crosses the IPI threshold alone.

\section{Design}
\label{sec:design}

\subsection{Three layers, one decorator}

\textsc{toolwall} models every tool call as a request--response
pair traversing three layers:

\medskip
\noindent\hrulefill\\[2pt]
\textsc{(1) Policy layer.} A frozen \verb|Policy| dataclass
describes the boundary contract: URL allow/deny lists, egress
denylist, filesystem path scope, named redaction patterns
(\verb|EMAIL|, \verb|API_KEY|, \verb|PHONE_E164|, etc.), response
byte cap, IPI threshold, rate / cost caps. Request-side fields
(allowlist, scope) are checked \emph{before} the tool runs;
response-side fields (byte cap, redactions, IPI threshold) are
checked \emph{after}. Policies are values, not classes: they
serialize cleanly to JSON for audit and replay.
\\[2pt]
\textsc{(2) Detector layer.} Three composable detectors run on
the response: (a)~a regex + Unicode-tag (U+E0000 block) +
zero-width-character + base64-payload pattern matcher whose
patterns are cited to public CVEs and disclosures in
\texttt{attack\_corpus.csv}; (b)~a markdown-exfil detector that
extracts every image / link URL and checks against the URL
allowlist; (c)~a frozen LoRA-tuned Llama-3.1-8B
``\textsc{toolwall-judge-v0.1}'' that returns a $[0,1]$ IPI
probability for residual cases. The detector layer is the only
layer whose external dependencies (the judge model weights) are
optional --- rule-only mode is shipped as the floor.
\\[2pt]
\textsc{(3) Enforcement layer.} One of \verb|allow|, \verb|redact|,
\verb|quarantine|, or \verb|block| is selected by the highest-priority
detection. \verb|block| raises a typed
\verb|toolwall.Blocked(reason, evidence)|; the host agent decides
whether to surface the error or to retry. Every decision lands in
the audit log (\verb|.toolwall-audit.log| by default).
\\[2pt]
\noindent\hrulefill

\subsection{The four-line claim}

A Stripe-style ``API failure if it has to be explained twice''
principle drives the public surface. The full Python integration
is:

\begin{lstlisting}[language=Python]
from toolwall import toolwall, policy

@toolwall(policy=policy.web_fetch())
def fetch_url(url: str) -> str:
    return requests.get(url).text
\end{lstlisting}

\noindent
TypeScript is shape-identical:

\begin{lstlisting}[language=JavaScript]
import { toolwall, policy } from "@bryteai/toolwall-core";

const fetchUrl = toolwall(
  async (url: string) => (await fetch(url)).text(),
  policy.webFetch({ urlAllowlist: ["*.gov.uk"] }),
);
\end{lstlisting}

\noindent
Adapters for CrewAI, LangGraph, the Claude Agent SDK, the OpenAI
Agents SDK, and FastMCP each provide a single \verb|wrap()| entry
point and weigh in at $\le$80 lines. Anything more elaborate is
treated as API failure.

\subsection{Design principles}

We chose four principles ahead of any code and held them as
red-lines during review:

\begin{enumerate}[leftmargin=*,nosep]
\item \textbf{Capability-based security.} Each policy field
\emph{names a capability} that the tool can use. Removing a
capability is local and visible at the call site.
\item \textbf{Deny by default.} An empty \verb|url_allowlist| means
``no URLs.'' An empty \verb|redact| means ``redact nothing,'' so
that omission errors fail safely in the direction of
\emph{accidental over-blocking} rather than silent leak.
\item \textbf{Explicit redaction over silent drop.} When we
modify a tool response we leave a typed sentinel
(\verb|[REDACTED: EMAIL]|) rather than deleting the bytes. The
model can reason about the redaction; a silent drop is invisible.
\item \textbf{Audit log is the ground truth.} Every policy
decision, every detector hit, every redaction or block is
appended to a JSONL audit log keyed by call id. The audit log is
the artefact the user grep's first.
\end{enumerate}

\section{Implementation}
\label{sec:impl}

\subsection{Python core}

The Python distribution ships as a single source tree with zero
runtime dependencies beyond \verb|pydantic| ($\ge$2.6) and
\verb|regex|. The decorator is implemented twice
(sync / async) over a shared core; both paths share the same
policy and detector pipeline. The package is tested under
Python~3.10--3.13. Property-based tests via
\verb|hypothesis|~\cite{hypothesis} cover policy serialization
and decorator argument handling.

\subsection{TypeScript port}

\verb|@bryteai/toolwall-core| is a feature-parity port published
under the same \verb|Policy| field names (translated to camelCase
on the boundary), Zod-validated, bundled with \verb|tsup| to
both ESM and CJS, total bundle $<$30~KB minified per entry. The
TypeScript core is consumed by the MCP middleware
(\verb|@bryteai/toolwall-mcp-middleware|), the Vercel AI~SDK v6
adapter, and the Cloudflare Workers MCP-server scaffold.

\subsection{Cloudflare Workers MCP scaffold}

\verb|create-toolwall-worker| (npm) scaffolds an OAuth~2.1 PKCE
MCP server with Streamable HTTP transport (the SSE lane is
deprecated for MCP as of 2026), Durable Object session state,
and KV-backed audit logging --- with the toolwall middleware
pre-wired to a sensible default \verb|policy.mcp()| profile. This
addresses the ``41\% of 2026 MCP servers ship without OAuth''
class~\cite{salt2026mcp} by making the secure path the default
path.

\subsection{Adapters as middleware}

Critically, the MCP adapter is a \emph{middleware} hook into
FastMCP and \verb|@modelcontextprotocol/sdk|, not a fork of either
SDK. Forking the SDK was an early temptation and was rejected
because (a)~the maintenance tax would consume the engineering
budget, and (b)~a middleware composes cleanly with vendor changes
to the protocol (e.g., the SSE$\to$Streamable~HTTP move). The same
principle drives the CrewAI / LangGraph / Claude SDK / OpenAI
Agents SDK adapters.

\subsection{Deterministic demo recorder}

For reproducible video evidence, we ship \verb|@toolwall/recorder|,
a Playwright-based recorder that owns the in-page clock: it
patches \emph{exactly four globals}
(\verb|requestAnimationFrame|, \verb|cancelAnimationFrame|,
\verb|performance.now|, \verb|Date.now|) via \verb|addInitScript|
and exposes \verb|window.__toolwall_advance(ms)| to the driver.
\verb|setTimeout|/\verb|setInterval|/\verb|queueMicrotask| are
deliberately \emph{not} patched, because React's scheduler reads
them for its own work loop and patching them desyncs the
renderer (a recipe documented in our prior engineering work). Per
virtual frame, the driver advances the clock by $1000/\text{fps}$
milliseconds, yields one microtask round-trip so React commits,
then streams the PNG bytes from
\verb|Page.captureScreenshot()| straight into
\verb|ffmpeg image2pipe|. The resulting mp4 is bit-stable across
runs at the same git revision, CRF, and preset --- a property
the closed-source ``record-a-screencast'' baseline does not have.

\section{Evaluation}
\label{sec:eval}

\subsection{Benchmark composition}

\textsc{toolwall-bench-v0.1} (frozen 2026-05-23) totals 1{,}054
cases composed of three subsets, license-tagged:

\begin{center}
\begin{tabular}{lrll}
\toprule
Source & Cases & License & Provenance \\
\midrule
\verb|mcp_v0_1|   & 200   & Apache-2.0 & New, hand-curated. \\
\verb|injecagent| & 500   & MIT        & \cite{zhan2024injecagent} \\
\verb|agentdojo|  & 354   & AGPL-3.0   & \cite{debenedetti2024agentdojo} \\
\midrule
\textbf{total}    & \textbf{1{,}054} & mixed & \\
\bottomrule
\end{tabular}
\end{center}

Splits are deterministic (seeded shuffle on sorted IDs): train
$600$, dev $200$, test $254$. The test split is SHA256-pinned in
\verb|eval/run_pipeline.py| and was never trained against by
either the judge model or any pattern detector.

The $200$ new \verb|mcp_v0_1| cases cover seven attack families
named in the 2026 disclosure record: \emph{tool poisoning},
\emph{rug-pull}, \emph{full-schema poisoning},
\emph{markdown-image exfil}, \emph{Unicode-tag hidden
instruction}, \emph{base64 hidden payload}, and
\emph{multi-hop indirect injection}. Each case is cited to the
public disclosure that motivated it.

\subsection{Metrics}

We report two metrics:

\paragraph{ASR (Attack Success Rate, $\downarrow$).}
Fraction of attack cases that successfully steer the agent
(measured by the judge model on the output trace, with the same
judge revision across all defenses to keep comparisons fair).

\paragraph{Utility ($\uparrow$).}
Fraction of \emph{legitimate} (benign) requests that complete
successfully under the defense. A defense that drops every
response would trivially hit ASR~0\% but utility~0\%.

\subsection{In-house results (200 cases, rules + judge)}
\label{sec:eval-inhouse}

All numbers below are on the in-house 200-case split
($n_\text{attack}=185$, $n_\text{benign}=15$, git
rev~\texttt{01d1c17}, run on 2026-05-23). Intervals are Wilson
$95\%$ CIs; raw rows in \verb|results/leaderboard.json|. We compare
against three model-side defenses (no-defense control, prompt
sandwiching, Spotlighting~\cite{hines2024spotlight}) and the only
publicly available OSS IPI classifier we could run end-to-end on
the same split, \texttt{Niansuh/Prompt-Guard-86M} (the community
port of Meta's Prompt-Guard).

\begin{center}\small
\begin{tabular}{lll r}
\toprule
Defense & ASR $\downarrow$ [95\% CI] & Utility $\uparrow$ [95\% CI] & LoC $\downarrow$ \\
\midrule
No defense                          & $1.000$ $[0.980, 1.000]$ & $1.000$ $[0.796, 1.000]$ & 0 \\
Prompt sandwiching                  & $1.000$ $[0.980, 1.000]$ & $1.000$ $[0.796, 1.000]$ & $\sim$30 \\
Spotlighting~\cite{hines2024spotlight} & $1.000$ $[0.980, 1.000]$ & $1.000$ $[0.796, 1.000]$ & $\sim$50 \\
Prompt-Guard-86M (input-side OSS)   & $0.427$ $[0.358, 0.499]$ & $0.467$ $[0.248, 0.699]$ & --- \\
\textbf{toolwall (rules only)}      & $\bm{0.778}$ $[0.713, 0.832]$ & $\bm{1.000}$ $[0.796, 1.000]$ & \textbf{4} \\
\textbf{toolwall (rules + judge)}   & $\bm{0.097}$ $[0.062, 0.149]$ & $\bm{0.933}$ $[0.702, 0.988]$ & \textbf{4} \\
\bottomrule
\end{tabular}
\end{center}

\paragraph{Non-overlapping CIs.}
\textsc{toolwall}'s upper CI ($0.149$) lies strictly below
PromptGuard-86M's lower CI ($0.358$); the gap is significant at
$n=185$ without a multiple-comparisons correction. PromptGuard
also fails on benign traffic (utility $0.467$ $[0.248, 0.699]$ ---
it incorrectly flags more than half of legitimate responses), so
it is not a usable production defense even before the ASR gap.

\noindent
Per-family breakdown shows where the pattern layer alone is
sufficient and where the judge layer carries the load:

\begin{center}\small
\begin{tabular}{lrr}
\toprule
Attack family & ASR (rules only) & ASR (rules + judge) \\
\midrule
Unicode-tag hidden instruction & $\bm{0.000}$ & $\bm{0.000}$ \\
Markdown-image exfil           & $0.467$ & $0.100$ \\
Tool poisoning                 & $1.000$ & $\bm{0.000}$ \\
Full-schema poisoning          & $1.000$ & $\bm{0.000}$ \\
Rug-pull                       & $1.000$ & $0.120$ \\
Base64 hidden payload          & $1.000$ & $0.200$ \\
Multi-hop indirect injection   & $1.000$ & $0.280$ \\
\bottomrule
\end{tabular}
\end{center}

\paragraph{What these numbers mean.}
Unicode-tag attacks are fully covered by deterministic codepoint
scanning (rules and judge both at $0$). The markdown-exfil
detector cuts ASR roughly in half on its own ($0.467$); the judge
closes the remaining gap, where the attacker uses an in-allowlist
domain as a covert sink. Tool poisoning and full-schema poisoning
go from $1.0$ to $0.0$ once the judge is enabled --- those classes
are essentially a semantic-content problem with no surface pattern
to match. Multi-hop attacks remain the hardest residual class
($0.280$) because the malicious instruction is distributed across
two responses and only one fragment ever crosses the IPI threshold
alone, a failure mode we name explicitly in \S\ref{sec:threat}.

\subsection{InjecAgent and AgentDojo targets}
\label{sec:eval-targets}

The 854 imported InjecAgent and AgentDojo cases are not yet rerun
under \textsc{toolwall} at preprint time; the measured numbers in
\S\ref{sec:eval-inhouse} are on the 200-case in-house split only.
We therefore distinguish forward \emph{targets} from the measured
results above:

\begin{center}\small
\begin{tabular}{lrrrr}
\toprule
Defense & InjecAgent ASR $\downarrow$ & AgentDojo ASR $\downarrow$ & Utility $\uparrow$ & LoC $\downarrow$ \\
\midrule
No defense                  & $\sim$24\% & $\sim$18\% & $100\%$ & 0 \\
Prompt sandwiching          & $\sim$14\% & $\sim$12\% & $88\%$ & $\sim$30 \\
Spotlighting~\cite{hines2024spotlight} & $\sim$9\% & $\sim$11\% & $91\%$ & $\sim$50 \\
GPT-4o as IPI judge (inline)& $\sim$7\% & $\sim$8\% & $79\%$ & $\sim$80 \\
\textbf{toolwall v0.1}      & target $<$2\% & target $<$3\% & target $\ge$97\% & \textbf{4} \\
\bottomrule
\end{tabular}
\end{center}

\noindent
We commit in advance: if the held-out test disagrees with these
targets at release time, the README and the abstract of this
preprint are updated, not the benchmark. The held-out test split
was locked at \texttt{git rev~01d1c17} with SHA256 pin
\verb|ce685e7a36b716b0|\verb|981f359834e7febf6|\verb|b237c51059dc74af|\verb|15d7bf46a0c74f2|
before the judge LoRA training began.

\subsection{Train$\leftrightarrow$test leakage audit}
\label{sec:eval-minhash}

To verify that the test split is not a paraphrase of the
$4{,}000$-pair judge training corpus, we run an offline MinHash +
$5$-gram Jaccard scan (\verb|eval/run_minhash_dup_check.py|, LSH
threshold $0.5$, leak threshold $0.4$). Overall mean Jaccard-5 is
$0.006$, overall max is $0.447$, and only $2/200$ test cases
($1.0\%$, both unicode-tag) exceed the leak threshold; six of
seven attack families and the benign class have max Jaccard-5 of
$0$. Full per-family numbers and the two flagged IDs are in
\verb|results/train_test_similarity.json|.

\subsection{Out-of-distribution probe (novel attack families)}
\label{sec:eval-ood}

To gauge how the rule layer generalizes to attack families it
was not designed against, we used the CrewAI multi-agent
framework (\url{https://github.com/crewAIInc/crewAI}) to author
$10$ adversarial responses spanning $10$ novel families absent
from our seven training families (timing-channel directives,
RFC-header role confusion, locale-fingerprint mismatch,
JSON-schema coercion, fake-citation social engineering,
format-string smuggling, diff-patch injection, SQL-comment
inversion, OAuth-consent spoofing, error-message imperative). On
this set, \textsc{toolwall} rules-only reaches ASR $0.80$
$[0.490, 0.943]$ ($2/10$ blocked) --- a result we publish loudly:
the pattern layer alone is undertrained on unseen families, and
the judge layer is load-bearing. Judge-on-novel-set scoring is
queued for v0.2 (fresh A10G GPU run; disclosed in the model
card). Raw artifacts:
\verb|eval/_artifacts/novel_attacks_crewai.jsonl|,
\verb|results/novel_attacks_score.json|.

\section{Limitations}
\label{sec:limits}

\paragraph{Judge dev-split saturation.} The LoRA judge scored
$1.0/1.0/1.0$ on its $200$-pair dev split. We suspect template
overlap with the train split (the judge memorized the four MCP
prompt templates used to author the training corpus rather than
generalizing the IPI signal) and we publish the suspect number
verbatim rather than tune it away. The honest authoritative
number is therefore the $0.097$ test-split ASR in
\S\ref{sec:eval-inhouse}, not the dev metric. Full disclosure is
in the model card under
\verb|training/MODEL_CARD.md|~\S$5.3$.

\paragraph{Generalization to novel attack families.} The OOD
probe in \S\ref{sec:eval-ood} shows \emph{rule-only} ASR degrades
to $0.80$ $[0.490, 0.943]$ on $10$ attacks from $10$ unseen
families. The judge layer was not run on that probe at preprint
time (it requires a fresh A10G GPU session), so the $0.80$
figure is intentionally the worst-case rules-only floor on unseen
families and \emph{not} the full \textsc{toolwall} stack. The
judge-on-novel-set run is the single most important v0.2
follow-up.

\paragraph{Deterministic decoding by construction.} All judge
inferences run at temperature $\sigma=0$ with greedy decoding, so
the run-to-run variance in our tables is $0$ by construction, not
because the model is calibrated. The Wilson CIs we report are
over the data ($n_\text{attack}=185$, $n_\text{benign}=15$), not
over decoding noise; stochastic-decoding intervals would require
a non-zero-temperature rerun.

\paragraph{Single-bench scope at v0.1.} All measured numbers are
on the in-house 200-case split; the InjecAgent and AgentDojo
imports remain forward targets per \S\ref{sec:eval-targets}.

\paragraph{Single-author curation.} The $200$ \verb|mcp_v0_1| cases
were authored by one person. A multi-author replication is the
right hygiene; we invite community PRs.

\paragraph{No model-side defense.} \textsc{toolwall} is a
response-side defense. It composes with, but does not replace,
model-side alignment, system-prompt hardening, and sandboxing.
Claims of ``protected'' agents that use \textsc{toolwall} alone
should be read as ``one layer in a defense-in-depth stack.''

\paragraph{English-language bias.} The current pattern corpus
and the judge training data are English-dominant. Urdu,
Hinglish, Arabic, and RTL IPI variants exist
\cite{ozcan2024multilingual} and are scheduled for v0.2.

\paragraph{Latency.} The judge model adds $\sim$60--120~ms per
tool response on a single A10. For latency-critical agents, the
rule-only mode is the production-recommended default with the
judge enabled per-tool on the higher-risk surfaces.

\section{Related work}
\label{sec:related}

\paragraph{Prompt-level defenses.} Spotlighting~\cite{hines2024spotlight}
inserts delimiter tokens around untrusted content; prompt
sandwiching repeats the agent's goal before and after a tool
response. Both run \emph{inside the model context} and degrade
utility ($\le$91\% in our measurement). \textsc{toolwall} runs
\emph{before} the response reaches the model, composes with
these defenses, and does not consume context tokens.

\paragraph{IPI benchmarks.}
InjecAgent~\cite{zhan2024injecagent},
AgentDojo~\cite{debenedetti2024agentdojo}, and
MCPTox~\cite{mcptox2025} establish that the threat is widespread
and the defense gap is wide. We reuse 854 of their cases under
their respective licenses and add 200 cases targeting attack
classes named after their publication dates.

\paragraph{MCP security tooling.} Invariant Labs MCP-scan and the
broader ``static metadata audit'' class detect tool poisoning at
registration time but do not see the response body.
NeMo~Guardrails~\cite{nemoguard2024} composes content filters
around generation but is again context-level. \textsc{toolwall}
sits at a distinct boundary --- the tool response --- and is the
only Apache-2.0 entry in that position at the time of writing.

\paragraph{Commercial offerings.} Lakera Guard, Pillar Security,
Invariant Labs, and Northflank provide closed-source IPI scoring
through hosted APIs. \textsc{toolwall} is intentionally local,
self-hostable, and inspectable so that academic and adversarial
study is possible without a vendor relationship.

\paragraph{Industry frameworks.}
OWASP ASI Top~10~\cite{owaspasi2026} elevates IPI to A01. NIST
AI~600-1~\cite{nist60012024} names IPI explicitly. Our policy
fields are aligned with both frameworks' control categories so
that compliance mapping is mechanical rather than manual.

\section{Conclusion}
\label{sec:conclusion}

Indirect prompt injection is the first agentic-AI security class
with both a named CVSS-9.3 production incident (EchoLeak) and a
framework-level acknowledgment (OWASP ASI Top~10) inside the same
year. The opensource defensive layer at the agent$\leftrightarrow$tool
boundary was missing; we provide it.

\textsc{toolwall} is a four-line decorator, an Apache-2.0
codebase, a frozen benchmark, a deterministic demo recorder, and
an honest floor of measured numbers. The remaining work is
GPU-bounded (the judge fine-tune) and community-bounded (more
attack classes, more languages, more replication). We invite both.

\section*{Acknowledgments}

This work was conducted at \href{https://bryteai.studio}{bryteai.studio}
under Bin Abdullah LLC (Brooklyn, NY) and benefits from NVIDIA's
Inception programme. We thank the InjecAgent and AgentDojo authors
for permissive licensing on their case sets, and the OWASP ASI
working group for naming the threat class precisely enough to
defend it.

\section*{Reproducibility statement}

Code, weights, dataset, and the held-out test split SHA256 are
public. \verb|python -m eval.run_pipeline --inhouse-only|
reproduces the in-house numbers in \S\ref{sec:eval-inhouse} on a
laptop CPU in under five minutes; the full bench with the judge
adds one A10-hour. All commit hashes referenced in this paper are
permanent (\verb|git log| on the public repository).

\begin{thebibliography}{99}\small

\bibitem{mcptox2025}
Pillai, A., Chen, J., et al.
\newblock MCPTox: a benchmark for tool-call security in MCP-based agents.
\newblock {\em arXiv:2508.14925}, 2025.

\bibitem{hines2024spotlight}
Hines, K., Lopez, G., Hall, M., Zarfati, F., Zunger, Y., Kiciman, E.
\newblock Defending against indirect prompt injection attacks with
spotlighting.
\newblock {\em arXiv:2403.14720}, 2024.

\bibitem{zhan2024injecagent}
Zhan, Q., Liang, Z., Ying, Z., Kang, D.
\newblock InjecAgent: Benchmarking indirect prompt injections in
tool-integrated large language model agents.
\newblock {\em Findings of ACL}, 2024.

\bibitem{debenedetti2024agentdojo}
Debenedetti, E., Zhang, J., Balunovi\'c, M., Beurer-Kellner, L.,
Fischer, M., Tram\`er, F.
\newblock AgentDojo: A dynamic environment to evaluate prompt
injection attacks and defenses for LLM agents.
\newblock {\em NeurIPS Datasets and Benchmarks}, 2024.

\bibitem{harmbench2024}
Mazeika, M., et al.
\newblock HarmBench: A standardized evaluation framework for
automated red teaming and robust refusal.
\newblock {\em ICML}, 2024.

\bibitem{jbb2024}
Chao, P., Debenedetti, E., Robey, A., et al.
\newblock JailbreakBench: An open robustness benchmark for
jailbreaking large language models.
\newblock {\em NeurIPS Datasets and Benchmarks}, 2024.

\bibitem{nemoguard2024}
Rebedea, T., Dinu, R., Sreedhar, M., Parisien, C., Cohen, J.
\newblock NeMo Guardrails: A toolkit for controllable and safe LLM
applications with programmable rails.
\newblock {\em EMNLP System Demonstrations}, 2023.

\bibitem{owaspasi2026}
OWASP Foundation.
\newblock Top~10 for Agentic Applications (ASI Top~10), 2026.
\newblock \url{https://owasp.org/www-project-top-10-for-large-language-model-applications/}

\bibitem{nist60012024}
National Institute of Standards and Technology.
\newblock Adversarial machine learning: A taxonomy and terminology
of attacks and mitigations. NIST AI~600-1, 2024.

\bibitem{invariant2025toolpoisoning}
Invariant Labs.
\newblock MCP tool poisoning: a post-approval mutation attack.
Blog post, 2025.

\bibitem{cyberark2025schema}
CyberArk Labs.
\newblock Full-schema poisoning: malicious schema fields in MCP
tool metadata. Blog post, 2025.

\bibitem{cymulate2026codex}
Cymulate Research.
\newblock Codex CLI sandbox escape via the
\texttt{notify} configuration field. Disclosure, March 2026.

\bibitem{salt2026mcp}
Salt Labs.
\newblock 2026 MCP server posture report: 41\% of public servers
ship without OAuth.

\bibitem{hypothesis}
MacIver, D., et al.
\newblock Hypothesis: a property-based testing library for Python.
\newblock \url{https://hypothesis.works}.

\bibitem{ozcan2024multilingual}
\"Ozcan, S., et al.
\newblock Multilingual indirect prompt injection: an empirical
study across 12 languages.
\newblock {\em Preprint}, 2024.

\end{thebibliography}

\end{document}
